\section{Architecture and training details}
\label{app:architectures}

Table~\ref{tab:architecture-settings} records the architecture settings in the current experiment manifest. Widths and internal computation differ; a shared number of training epochs does not imply equal parameter counts or compute. The baseline named GraphSAGE is the weighted implementation described in Section~\ref{sec:architectures}, and \method{} denotes the combined extension used in the study.

\begin{table}[h]
 \centering
 \caption{Fixed retrieval configurations. ``3/3/3'' denotes ReaRev's instructions, reasoning steps, and adaptive iterations, respectively. Dropout 0.1 is explicitly configured for the first five models; NBFNet uses its fixed path-propagation implementation. Settings are reproduced from the experiment manifest.}
 \label{tab:architecture-settings}
 \begingroup
\small
\begin{tabularx}{\textwidth}{@{}lccX@{}}
\toprule
Architecture & Width & Depth & Main configured characteristics \\
\midrule
GraphSAGE & 256 & 3 & Fixed cosine edge weights; MLP node classifier. \\
\method{} & 256 & 3 & Learned edge MLP (width 256), reverse edges, residual normalization, question-aware classifier. \\
R-GCN & 256 & 3 & 30 relation bases; categorical forward/reverse relations. \\
HGT & 128 & $3^{\dagger}$ & 8 attention heads; one entity node type; relation-dependent attention. \\
ReaRev & 50 & 3/3/3 & Frozen MiniLM; token instructions; adaptive revision. \\
NBFNet & 32 & 3 & Question-conditioned boundary states; DistMult messages; PNA aggregation. \\
\bottomrule
\end{tabularx}
\endgroup

 \par\smallskip
 {\footnotesize\raggedright $\dagger$ The current manifest uses three HGT layers. Its seed-1337 entry was corrected from two to three after the recorded run date; the saved summary does not establish that run's depth. Historical configuration verification is pending, so the HGT variability must not yet be interpreted as purely seed-induced.\par}
\end{table}

\paragraph{Aggregation variants.}
For an incoming triple $(u,r,v)$, the GraphSAGE implementation sums weighted transformed source states and concatenates this sum with the current target state before a learned update. The baseline weight is cosine question--relation similarity. \method{} uses Equation~\ref{eq:gate} and the residual-normalized update $\mathbf h_v^{(\ell+1)}=\operatorname{ReLU}(\operatorname{LayerNorm}(\mathbf h_v^{(\ell)}+\operatorname{Dropout}(\mathbf u_v^{(\ell)})))$, where $\mathbf u_v^{(\ell)}$ is the message-passing update. R-GCN uses $W_r^{(\ell)}=\sum_{b=1}^{30}a_{rb}^{(\ell)}V_b^{(\ell)}$ and relation-wise neighbor normalization. Its forward and reverse relations have separate identifiers. HGT retains relation-specific attention and message transformations with eight heads, shared node-type projections, a learned residual gate, and normalization.

\paragraph{Task-oriented models.}
ReaRev uses the frozen encoder \texttt{sentence-transformers/all-MiniLM-L6-v2}, revision \texttt{1110a243fdf4706b3f48f1d95db1a4f5529b4d41}, with maximum question and relation lengths of 64 and 16 tokens. Token-level attention produces three instructions. Initial entity states derive from incident relation information, and reasoning begins from a seed distribution. Instructions are revised after each reasoning iteration. NBFNet supplies the full projected question vector to each seed, uses separate original and reverse relation operators, and reinjects the initial boundary condition at each layer. PNA combines mean, maximum, minimum, and standard deviation statistics with degree scaling \citep{corso2020pna}. Final node states are concatenated with the question representation for scoring.

\paragraph{Training objective.}
All models except ReaRev use independent binary answer labels and class-balanced binary cross-entropy. For a graph with $n_+$ positive and $n_-$ negative nodes, the positive term receives weight $n_-/n_+$ when $n_+>0$. ReaRev instead minimizes KL divergence to a distribution uniform over the gold answer nodes. Independent outputs are passed through sigmoid, whereas ReaRev yields a softmax distribution. The common candidate-selection rule is applied after these normalizations. The model comparison therefore includes differences in objective and calibration, as well as architecture.

\section{Metric definitions and evaluation protocol}
\label{app:metrics}

Let $Q$ be the set of evaluated questions with nonempty normalized gold sets and $N=|Q|=1746$. Let $C_q^{(k)}$ be the top $k$ retrieved candidates, $C_q$ the final candidate set, $U_q$ the source and target entities actually appearing in emitted evidence triples, and $P_q$ the normalized prediction set. The indicator function $\ind[\cdot]$ is one when its argument holds. Reported standard deviations are computed over the three run-level values, using denominator $3-1$.

\paragraph{Retrieval.}
Hits@k measures whether at least one gold answer appears in the first $k$ candidates:
\begin{equation}
 \operatorname{Hits@}k=\frac1N\sum_{q\in Q}\ind[A_q\cap C_q^{(k)}\ne\varnothing],
 \quad
 \operatorname{nDCG@}k=\frac1N\sum_{q\in Q}
 \frac{\sum_{i=1}^{k}\mathrm{rel}_{qi}/\log_2(i+1)}
 {\sum_{i=1}^{\min(|A_q|,k)}1/\log_2(i+1)}.
\end{equation}
Relevance is binary, and nonexistent candidate positions contribute zero. nDCG rewards earlier relevant results and normalizes by an ideal list \citep{jarvelin2002cumulated}. The coverage metrics in Table~\ref{tab:retrieval} are
\begin{equation}
 \operatorname{RGC}=\frac1N\sum_{q\in Q}\frac{|A_q\cap C_q|}{|A_q|},
 \qquad
 \operatorname{RFGC}=\frac1N\sum_{q\in Q}\ind[A_q\subseteq C_q].
\end{equation}
Both are macro averages: each question has equal weight regardless of its answer-set size.

\paragraph{Evidence.}
AverageTriples is $N^{-1}\sum_q|H_q|$ and AverageDistinctNodes is $N^{-1}\sum_q|U_q|$. Let $D_q$ be the retrieved candidates retained according to the construction output. ContextCandidateCoverage is the macro average of $|D_q|/|C_q|$, with zero for an empty candidate list. CandidateReduction instead pools counts: $100\sum_q(|C_q|-|D_q|)/\sum_q|C_q|$. These two aggregates need not be exact complements. ContextGoldCoverage (CGC) and ContextFullGoldCoverage (CFGC) replace $C_q$ by $U_q$ in the RGC and RFGC formulas. Candidate retention and evidence visibility are distinct: a seed-only, zero-edge structure does not expose an entity in a serialized triple.

\paragraph{Final answer quality.}
Hit is $N^{-1}\sum_q\ind[A_q\cap P_q\ne\varnothing]$, and Exact Match (EM) is $N^{-1}\sum_q\ind[A_q=P_q]$. For micro metrics, $\mathrm{TP}=\sum_q|A_q\cap P_q|$, $\mathrm{FP}=\sum_q|P_q\setminus A_q|$, and $\mathrm{FN}=\sum_q|A_q\setminus P_q|$. Precision is $\mathrm{TP}/(\mathrm{TP}+\mathrm{FP})$, recall is $\mathrm{TP}/(\mathrm{TP}+\mathrm{FN})$, and F1 is their harmonic mean. Zero denominators yield zero. Failed generations use an empty prediction set and are forced to fail EM. These pooled metrics weight individual answers rather than questions.

\paragraph{Context utilization.}
Define $F_q=\ind[A_q\subseteq U_q]$ and $O_q=\ind[A_q\nsubseteq P_q]$. The two full-context outcomes and the conditional omission rate are
\begin{align}
 \operatorname{FullContextCompleteAnswer}&=\frac1N\sum_q F_q(1-O_q),\\
 \operatorname{FullContextLLMOmission}&=\frac1N\sum_q F_qO_q,\\
 \operatorname{LLMOmissionGivenFullContext}&=\frac{\sum_q F_qO_q}{\sum_q F_q}.
 \label{eq:conditional-omission}
\end{align}
The last metric is undefined if no question has full context coverage; that case does not occur in the displayed results. The first two rates sum to CFGC for each run, not to one. Completeness allows extra predicted entities, while Exact Match does not. Conditional rates are computed within each run before averaging; dividing already averaged rates can give a different value.

\paragraph{Normalization and accounting.}
Answer strings are lowercased and trimmed; punctuation and English articles are removed, whitespace is collapsed, and duplicate normalized answers are removed. Comparisons use these normalized names rather than a learned semantic-equivalence judge. Per-stage records are joined by instance identifier. TotalTokens sums recorded input and output usage for a run, including instructions and question tokens as well as evidence. Resource comparisons are made primarily within a language model because tokenizers differ.

\section{Complete evidence-construction results}
\label{app:evidence-results}

Table~\ref{tab:evidence-full} reports all nine evidence settings. The three NBFNet seeds are reused throughout. PCST operates on a multigraph of original facts, with no synthetic inverse facts in the output. The implementation calls \texttt{pcst\_fast} with a specified root, one cluster, and \texttt{gw} pruning. For multiple seeds, each seed receives an additional prize greater than the sum of candidate prizes and structural edge costs, and is connected to a temporary root at zero cost. Removing that root can leave multiple seed-anchored branches. The quoted candidate prizes and reported collected prize exclude these enforcement additions.

\begin{table}[h]
 \centering
 \caption{Evidence results, mean $\pm$ standard deviation across three NBFNet seeds. CCC: macro candidate coverage; Red.: pooled candidate reduction in percent; CGC: context gold coverage; CFGC: full context gold coverage. Constant and semantic rows differ only in their edge-cost definition and scale.}
 \label{tab:evidence-full}
 \begingroup
\footnotesize
\setlength{\tabcolsep}{3pt}
\begin{tabular}{@{}lccccccc@{}}
\toprule
Strategy & $\lambda$ & Triples & Nodes & CCC & Red. (\%) & CGC & CFGC \\
\midrule
Shortest paths & -- & $94.47\pm22.58$ & $32.12\pm6.60$ & $1.000\pm0.000$ & $0.00\pm0.00$ & $0.888\pm0.006$ & $0.838\pm0.007$ \\
\midrule
PCST constant & 0.01 & $14.45\pm0.26$ & $15.46\pm0.26$ & $1.000\pm0.000$ & $0.00\pm0.00$ & $0.881\pm0.004$ & $0.831\pm0.005$ \\
 & 0.25 & $14.45\pm0.26$ & $15.46\pm0.26$ & $1.000\pm0.000$ & $0.00\pm0.00$ & $0.881\pm0.004$ & $0.831\pm0.005$ \\
 & 0.5 & $14.30\pm0.25$ & $15.31\pm0.25$ & $0.994\pm0.003$ & $0.61\pm0.30$ & $0.881\pm0.004$ & $0.830\pm0.006$ \\
 & 1 & $13.57\pm0.26$ & $14.58\pm0.26$ & $0.949\pm0.003$ & $5.02\pm0.33$ & $0.875\pm0.003$ & $0.816\pm0.005$ \\
\midrule
PCST semantic & 0.01 & $13.70\pm0.28$ & $14.71\pm0.28$ & $1.000\pm0.000$ & $0.00\pm0.00$ & $0.881\pm0.004$ & $0.831\pm0.006$ \\
 & 0.25 & $13.70\pm0.28$ & $14.71\pm0.28$ & $1.000\pm0.000$ & $0.00\pm0.00$ & $0.881\pm0.004$ & $0.831\pm0.006$ \\
 & 0.5 & $13.70\pm0.28$ & $14.71\pm0.28$ & $1.000\pm0.000$ & $0.00\pm0.00$ & $0.881\pm0.004$ & $0.831\pm0.006$ \\
 & 1 & $13.50\pm0.24$ & $14.51\pm0.24$ & $0.991\pm0.004$ & $0.82\pm0.38$ & $0.881\pm0.004$ & $0.830\pm0.006$ \\
\bottomrule
\end{tabular}
\endgroup

\end{table}

The table highlights why candidate survival alone cannot determine context usefulness. Both costs at $\lambda=0.01$ retain all retrieved candidates, and their CGC values round to the same number, while downstream answer quality differs substantially. The measured difference is associated with the selected connecting structure, not with a wholesale reduction in candidate availability. Similar rounded aggregates are not evidence that the individual subgraphs are identical.

\section{Complete final results and reproducibility details}
\label{app:final-results}

Table~\ref{tab:final-full} gives all ten model--evidence combinations. Each row summarizes three runs associated with the same three NBFNet retriever seeds. Table~\ref{tab:omission-full} provides the exact conditional rates used in Section~\ref{sec:context-results}, including the low-cost endpoints. No additional training or inference runs were introduced to prepare these tables.

\begin{table}[h]
 \centering
 \caption{Complete final-answer results. P and R are micro precision and recall; EM is Exact Match; tokens are total input plus output tokens in millions per run. SP denotes shortest-path union, C constant PCST, and S semantic PCST. All entries are mean $\pm$ standard deviation across three runs.}
 \label{tab:final-full}
 \begingroup
\footnotesize
\setlength{\tabcolsep}{4pt}
\begin{tabular}{@{}lcccccc@{}}
\toprule
Evidence & Hit & P & R & F1 & EM & Tokens (M) \\
\midrule
\multicolumn{7}{@{}l}{\textbf{DeepSeek-V4-Flash}} \\
SP & $0.802\pm0.005$ & $0.860\pm0.009$ & $0.493\pm0.007$ & $0.627\pm0.008$ & $0.580\pm0.003$ & $3.42\pm0.77$ \\
C, $\lambda=0.01$ & $0.598\pm0.012$ & $0.881\pm0.006$ & $0.259\pm0.005$ & $0.400\pm0.007$ & $0.322\pm0.002$ & $0.72\pm0.01$ \\
C, $\lambda=1$ & $0.594\pm0.014$ & $0.884\pm0.007$ & $0.257\pm0.006$ & $0.398\pm0.008$ & $0.320\pm0.002$ & $0.69\pm0.01$ \\
S, $\lambda=0.01$ & $0.688\pm0.007$ & $0.885\pm0.008$ & $0.342\pm0.008$ & $0.493\pm0.010$ & $0.410\pm0.003$ & $0.69\pm0.01$ \\
S, $\lambda=1$ & $0.689\pm0.006$ & $0.886\pm0.008$ & $0.341\pm0.008$ & $0.492\pm0.010$ & $0.411\pm0.004$ & $0.69\pm0.01$ \\
\midrule
\multicolumn{7}{@{}l}{\textbf{GPT-5.6 Luna}} \\
SP & $0.804\pm0.015$ & $0.846\pm0.007$ & $0.481\pm0.006$ & $0.613\pm0.006$ & $0.571\pm0.006$ & $3.07\pm0.66$ \\
C, $\lambda=0.01$ & $0.651\pm0.009$ & $0.847\pm0.001$ & $0.261\pm0.004$ & $0.399\pm0.005$ & $0.341\pm0.006$ & $0.70\pm0.01$ \\
C, $\lambda=1$ & $0.645\pm0.015$ & $0.849\pm0.006$ & $0.260\pm0.004$ & $0.398\pm0.006$ & $0.331\pm0.009$ & $0.67\pm0.01$ \\
S, $\lambda=0.01$ & $0.723\pm0.004$ & $0.850\pm0.008$ & $0.331\pm0.004$ & $0.476\pm0.005$ & $0.414\pm0.003$ & $0.67\pm0.01$ \\
S, $\lambda=1$ & $0.728\pm0.006$ & $0.854\pm0.007$ & $0.331\pm0.006$ & $0.477\pm0.006$ & $0.414\pm0.005$ & $0.66\pm0.01$ \\
\bottomrule
\end{tabular}
\endgroup

\end{table}

\begin{table}[h]
 \centering
 \caption{LLMOmissionGivenFullContext for every evaluated evidence setting. Lower is better. Conditioning is on full gold-entity coverage in each run's evidence, not on a shared fixed subset across all construction strategies.}
 \label{tab:omission-full}
 \begingroup
\small
\begin{tabular}{@{}lcc@{}}
\toprule
Evidence & DeepSeek-V4-Flash & GPT-5.6 Luna \\
\midrule
Shortest paths & $0.211\pm0.002$ & $0.237\pm0.010$ \\
PCST constant, $\lambda=0.01$ & $0.589\pm0.008$ & $0.568\pm0.012$ \\
PCST constant, $\lambda=1$ & $0.587\pm0.007$ & $0.573\pm0.014$ \\
PCST semantic, $\lambda=0.01$ & $0.477\pm0.001$ & $0.477\pm0.002$ \\
PCST semantic, $\lambda=1$ & $0.477\pm0.003$ & $0.474\pm0.005$ \\
\bottomrule
\end{tabular}
\endgroup

\end{table}

\paragraph{Generation protocol.}
The configured API model identifiers are \texttt{deepseek-v4-flash} and \texttt{gpt-5.6-luna}; these names identify the recorded runs. The generation service requests temperature zero. It supplies the question, deduplicated directed triples, and a common answer-only instruction requiring JSON of the form \texttt{\{"answers": ["complete entity name"]\}}. Entity names containing commas must remain a single item. Unsupported questions should return an empty array, and explanations are disabled. The system instruction reads:
\begin{quote}\small
You answer questions using only the provided reasoning paths. Return only valid JSON with the key answers. The answers value must be a JSON array of complete entity names. Do not generate an explanation. If the paths do not support an answer, return an empty answers array.
\end{quote}

\paragraph{Recorded artifacts.}
The implementation separates training, retriever evaluation, evidence construction, and inference outputs. Experiment manifests pin the three seeds and map the six retrievers, nine evidence conditions, and ten final conditions to persisted runs. The analysis comprises 18 retrieval runs, 27 evidence runs, and 30 final-generation runs. Result exports retain per-run values, aggregated summaries, and source-run provenance, allowing tables to be regenerated from existing measurements. Numeric experiment identifiers in the implementation are zero-based; the scientific comparisons here are organized by stage rather than identifier.

\paragraph{Scope of reproducibility.}
Training uses PyTorch and the settings in Section~\ref{sec:setup}; graph preparation and name mapping are part of the recorded pipeline. Exact reproduction also depends on the original prepared dataset, encoder artifacts, and API model availability. The comparison is between fixed implementations, with their different input encoders and objectives, and does not establish equal computational budgets. Public repository and tracking links are omitted from this anonymous draft. A review-ready anonymous artifact link must be supplied separately before submission; no anonymous archive is claimed to have been deposited here.
